\SourceLeaf{248}%
\Needspace{7\baselineskip}
\begin{center}
[\textsc{Gerling} \emph{to} \textsc{Gauss}.\quad \emph{Marburg, 26 February 1844.}]
\end{center}

\{\mbox{.\ .\ .\ .\ .\ .\ .\ .\ .} Ever since I came upon the symmetric triangles and the area of the spherical triangle in 1813, I have always felt the want of a precise definition or a precise criterion of symmetry as distinguished from congruence. Since last winter I have hit upon the following way of presenting the matter in my \textls{lectures} as well, for I had already arrived at it for myself on another occasion (1834).

We refer every figure in the plane to \textls{two} dimensions (rectangular coordinates) and every solid figure to \textls{three}. If we now suppose that the coordinates are subject to changes of sign, a figure that changes sign in \textls{two} dimensions generally becomes congruent, whereas a figure that changes sign in \textls{one} or in \textls{three} dimensions becomes symmetric. --- It follows as a corollary

1) that for coordinates taken from one point, reversing all the radii vectores (sit venia verbo) must yield symmetry in solid constructions and congruence in plane constructions,

\Needspace{7\baselineskip}
2) that symmetry passes into congruence when the figure can be divided by a straight line or a plane into two symmetric halves. ---

I should be much obliged if you would let me know on some occasion what you think of this way of presenting the matter.\}

\par\bigskip

\Needspace{7\baselineskip}
\begin{center}
\textsc{Gauss} to \textsc{Gerling}.\quad Göttingen, 8 April 1844.
\end{center}

\mbox{.\ .\ .\ .\ .\ .\ .\ .} Your remarks on symmetry and congruence are entirely to the point. What remains to be desired is the metaphysical reason \textls{why} it is so (which appears in your account merely as an observed fact), and with it the extension to a geometry of more than 3 dimensions, of which we human beings have no intuition, but which,\space%
